\section{Classic Block-Encoding Leaves}
\label{sec:classic-be-leaves}

This note records the paper-facing versions of the reusable ABEIS textbook
leaves.  The matching Lean declarations live mainly in
\texttt{QuantumBlockEncoding/BlockEncodingClassics.lean}; the matching agent
memory cards live under
\texttt{research-wiki/block-encoding-library/cards/}.  These statements are
intended as proof templates.  A concrete task must still instantiate the
registers, normalizer, cleanup condition, and resource tuple.

\subsection{Entrywise clean-block proof}

Let \(U\) act on an ancilla register \(A_{\rm anc}\) and a system register
\(S\).  Fix a clean ancilla basis state \(0_a\).  Define
\[
  \operatorname{Clean}_{0_a}(U)_{ij}
  := U_{(0_a,i),(0_a,j)} .
\]
If
\[
  U_{(0_a,i),(0_a,j)} = A_{ij}/\alpha
  \qquad\text{for all }i,j,
\]
then \(\operatorname{Clean}_{0_a}(U)=A/\alpha\).  Therefore, once unitarity of
\(U\) is also proved, \(U\) is an exact \((\alpha,a,0)\)-block encoding of
\(A\).

The proof is just matrix extensionality: two finite matrices are equal when
all entries agree.  In Lean this is the first reduction that should be reused
instead of re-expanding a circuit at every theorem.

\subsection{Partial permutation matrix unit}

Suppose the desired operator is
\[
  E = \ket{x_{\rm dst}}\!\bra{x_{\rm src}}\otimes I_S .
\]
Choose a permutation \(p\) of the full ancilla--active--passive basis such that
\[
  p(0_a,x_{\rm src},s)=(0_a,x_{\rm dst},s)
  \qquad\text{for every }s,
\]
and such that \(p(0_a,x,s)\) has non-clean ancilla whenever
\(x\neq x_{\rm src}\).  Let \(U_p\) be the permutation matrix of \(p\).  Then
\[
  (U_p)_{(0_a,x',s'),(0_a,x,s)}
  =
  \begin{cases}
  1, & x=x_{\rm src},\ x'=x_{\rm dst},\ s'=s,\\
  0, & \text{otherwise}.
  \end{cases}
\]
Thus the clean block is exactly \(E\).  This is the main proof pattern for the
transfer-operator main case.

\subsection{One-sparse support}

Let each column \(j\) have at most one possible nonzero row \(c(j)\).  Assume
\[
  A_{ij}=0\quad\text{whenever }i\neq c(j).
\]
An amplitude oracle supplies \(A_{c(j),j}\) on the clean signal branch, and a
support permutation maps \(j\) to \(c(j)\).  The clean entry becomes
\[
  A_{c(j),j}\,\delta_{i,c(j)} = A_{ij}.
\]
The reusable proof leaf is the Kronecker-delta support collapse.

\subsection{Sparse column oracle}

For an \(s\)-sparse column oracle, let \(c(j,\ell)\) list possible nonzero rows
for column \(j\).  Uniformly preparing the slot register produces a clean entry
of the form
\[
  \frac{1}{s}\sum_{\ell} A_{c(j,\ell),j}\,\delta_{i,c(j,\ell)}.
\]
If exactly one slot \(\ell\) satisfies \(c(j,\ell)=i\), this collapses to
\(A_{ij}/s\).  If no slot hits \(i\), it collapses to \(0\).  A concrete task
must prove the relevant support uniqueness condition.

\subsection{LCU prepare--select}

Let \(T=\sum_{\ell}\beta_\ell A_\ell\), and suppose each \(A_\ell\) already
has a compatible block encoding.  PREPARE creates amplitudes proportional to
\(\sqrt{\beta_\ell}\), SELECT applies the chosen block, and
PREPARE\({}^\dagger\) projects back.  The clean component is
\[
  \frac{1}{\sum_\ell \beta_\ell}
  \sum_\ell \beta_\ell A_\ell.
\]
The simplest compiled ABEIS leaf is the two-term weighted-sum entry identity.
Full PREPARE--SELECT--PREPARE semantics for arbitrary coefficient encodings is
kept as a larger theorem target.

\subsection{Contraction dilation}

If \(x^2+y^2=1\), the matrix
\[
  R(x,y)=
  \begin{pmatrix}
    x & y\\
    y & -x
  \end{pmatrix}
\]
has orthonormal rows, because each row has squared norm \(x^2+y^2=1\) and the
row inner product is \(xy-yx=0\).  Its clean entry is \(x\).  This scalar leaf
is the \(2\times2\) core of the SVD/contraction dilation fallback.  For a
general matrix contraction, the same idea is applied in a singular-value
basis; a task still needs the corresponding analytic or spectral backend.

\subsection{QSVT as a consumer}

QSVT is not a shortcut for constructing the original oracle.  It is a consumer
of a proved block encoding.  A task may cite a QSVT contract only after it has
named the input block encoding, the polynomial side conditions, and the
normalizer/error propagation.  Until those side conditions are formalized, the
QSVT node remains an external contract in the proof DAG.
